受格罗莫夫与埃利亚什伯格刚性定理 (El87) 的启发，自然可以把\emph{辛同胚}（简称 sympeomorphism）定义为：在同胚群内，任意辛微分同胚序列的
$C^0$-极限。在 (OhM07) 中，作者们表述了哈密顿微分同胚群的
$C^0$-对应物以及\emph{拓扑哈密顿流}的概念。事实表明，\emph{哈密顿同胚}（简称 hameomorphism）的集合构成辛同胚群的一个正规子群，并且人们相信它是辛同胚群的真子群。倘若真是这样，它将对动力系统中的一个重要问题给出否定答案，即 $S^2$（或 $D^2$）的保面积同胚群是否为单群。

这为辛拓扑提供了一个重要研究方向，人们希望它能引向
\stfhIndex{I085}{lian^xu^ha^mi^dun^dong^li^xue}{连续哈密顿动力学（continuous Hamiltonian dynamics）}
$C^0$-辛拓扑空间的正确概念。本章探索辛拓扑与拓扑哈密顿动力学这一片大体尚未开拓的领域。我们先紧随埃克兰与霍费尔 (EkH90) 的证明，说明一个 $C^0$-辛刚性定理的证明；然后说明穆勒与本书作者在 (OhM07) 和 (Oh10) 中引入的连续哈密顿流之定义。

\section{\texorpdfstring{$C^0$-辛刚性定理}{C0 symplectic rigidity theorem}}
\label{sec:c0-symplectic-rigidity}

我们先陈述著名的埃利亚什伯格--格罗莫夫 $C^0$-辛刚性定理。
\stfhIndex{I005}{00-symbol-005}{$C^0$-辛刚性定理（$C^0$ symplectic rigidity theorem）}
\stfhIndex{I112}{ai^li^ya^shen^bo^ge}{埃利亚什伯格（Eliashberg）}
\stfhIndex{I113}{ai^li^ya^shen^bo^ge^ge^luo^mo^fu}{埃利亚什伯格--格罗莫夫（Eliashberg–Gromov）}

\begin{theorem}[埃利亚什伯格--格罗莫夫]\label{thm:c0-symplectic-rigidity}
辛微分同胚子群 $\Symp(M,\omega)$ 在 $\Diff(M)$ 中是 $C^0$-闭的。更准确地说，如果辛微分同胚序列 $\phi_j$ 有一个可微的
$C^0$-极限 $\phi_\infty$，则 $\phi_\infty$ 是一个
$C^1$-辛微分同胚。
\end{theorem}

埃利亚什伯格的证明 (El87) 依赖于某些勒让德子流形波前的一个结构定理。这个结构定理的完整证明细节尚未见于文献。不过，格罗莫夫非挤压定理可以取代埃利亚什伯格的论证，给出另一个证明。事实上，只要在
$\mathbb C^n$ 的开集集合上存在\emph{任意辛容量}函数，就能利用埃利亚什伯格在
(El87) 中的论证给出一个相对直接的证明。

在本节余下部分，我们利用非挤压定理定义 $\mathbb R^{2n}$ 上的一个这种辛容量。关于从哈密顿动力学和经典临界点理论观点出发的辛刚性理论之精彩论述，请读者参阅 (HZ94)。

设 $F\subset(M,\omega)$ 是一个闭子集。

\begin{definition}[格罗莫夫面积]\label{def:gromov-area-closed-set}
\stfhIndex{I142}{ge^luo^mo^fu^mian^ji}{格罗莫夫面积（Gromov area）}
对 $M$ 的任意闭子集 $F$，定义
\[
 \underline c(F,\omega)
 =\sup\{\pi r^2\mid
   \text{存在辛嵌入 }\phi:(B^{2n}(r),\omega_0)\to(F,\omega)\},
\]
并称其为 $F$ 的\emph{格罗莫夫面积}。
\end{definition}

显然，如果 $F$ 的内部为空，则 $\underline c(F,\omega)=0$。

本节余下部分把这个定义限制于 $\mathbb R^{2n}$ 中的子集。上述定义的非平凡性是现代辛拓扑的基石之一，其证明是格罗莫夫著名非挤压定理的直接推论。关于该定理的证明，请读者参阅 (Gr85) 以及第 11.1 节。

接着定义另一个在刚性定理证明中更有用的容量：
\[
 \overline c(F)=\inf\{\pi r^2\mid
   \text{存在 }\psi\in\Symp(\mathbb R^{2n})
   \text{ 使得 }\psi(F)\subset Z^{2n}(r)\}.
\]
显然有 $\overline c(F)\geq\underline c(F)$。

\begin{proposition}\label{prop:extremal-capacity-properties}
设 $c$ 是 $\mathbb C^n$ 上的 $\underline c$ 与 $\overline c$ 中任意一个。则 $c$ 具有下列性质：
\begin{enumerate}[label=(\arabic*)]
 \item 单调性：若 $S\subset T$，则 $c(S)\leq c(T)$；
 \item 共形性：若 $\phi^*\omega_0=\alpha\omega_0$ 且
 $\alpha\ne0$，则 $c(\phi(F))=|\alpha|c(F)$；
 \item 归一化：
 $\underline c(B^{2n}(1),\omega_0)=\pi
 =\overline c(Z^{2n}(1),\omega_0)$。
\end{enumerate}
\end{proposition}

\begin{proof}
我们只给出 $\underline c$ 的证明，把另一个留给读者。单调性与共形性是 $\underline c$ 定义的简单推论。因此，只证明归一化公理。先注意，单调性公理蕴含
$\underline c(Z^{2n}(1),\omega_0)\geq
\underline c(B^{2n}(1),\omega_0)$。另一方面，对所有 $r>1$，非挤压定理蕴含
\[
 \underline c(Z^{2n}(1))
 <\underline c(B^{2n}(r))=\pi r^2.
\]
所以还有 $\underline c(Z^{2n}(1),\omega_0)\leq\pi$，从而已经证明
$\underline c(Z^{2n}(1),\omega_0)=\pi$。恒等式
$\underline c(B^{2n}(1),\omega_0)=\pi$ 是平凡的。
\end{proof}

\begin{definition}[\texorpdfstring{$\mathbb C^n$}{C-n} 中的辛容量]
\stfhIndex{I273}{xin^rong^liang}{辛容量（symplectic capacity）}
\label{def:symplectic-capacity-cn}
定义在 $\mathbb C^n$ 的子集上的任何函数 $c$，如果满足命题 6.1.3 中陈述的三个性质，就称为 $\mathbb C^n$ 上的\emph{辛容量函数}。
\end{definition}

命题 6.1.3 给出了 $\mathbb C^n$ 上辛容量的两个例子。下面的命题是定义的直接推论。

\begin{proposition}\label{prop:capacity-extremal-bounds}
对 $\mathbb C^n$ 上任意辛容量 $c$ 和所有子集
$F\subset\mathbb C^n$，都有
\[
 \underline c(F)\leq c(F)\leq\overline c(F).
\]
\end{proposition}

\begin{proof}
由归一化公理，对任意辛球容易证明
\begin{equation}
 \overline c(S)=c(S)=\underline c(S).
 \label{eq:capacities-agree-on-balls}
\end{equation}
设 $F\subset\mathbb C^n$ 是有界闭子集。由 $\overline c$ 的定义，存在辛微分同胚
$\psi_1:\mathbb C^n\to\mathbb C^n$，使得
\[
 \psi_1(F)\subset Z^{2n}(r+\epsilon),
\]
其中 $\pi r^2=\overline c(F)$。所以导出
\[
 c(F)=c(\psi_1(F))\leq\pi(r+\epsilon)^2.
\]
由于这对所有 $\epsilon>0$ 都成立，得到
$c(F)\leq\pi r^2=\overline c(F)$。

接着证明 $\underline c(F)\leq c(F)$。如果 $\underline c(F)=0$，则无须证明，故可假设 $\underline c(F)>0$。由 $\underline c(F)$ 的定义，对任意给定的
$0<\epsilon<\underline c(F)$，可以找到辛微分同胚 $\psi_2$，使得
\[
 \psi_2(B^{2n}(r-\epsilon))\subset F,
\]
其中 $\pi r^2=\underline c(F)$。由辛不变性与单调性，得到
\[
 c(B^{2n}(r-\epsilon))
 =c(\psi_2(B^{2n}(r-\epsilon)))\leq c(F).
\]
另一方面，$c$ 的归一化公理蕴含
\[
 c(B^{2n}(r-\epsilon))=\pi(r-\epsilon)^2,
\]
从而 $\pi(r-\epsilon)^2\leq c(F)$。由于这对任意
$0\leq\epsilon<r$ 都成立，得到 $\pi r^2\leq c(F)$，所以
$\underline c(F)\leq c(F)$。证明完毕。
\end{proof}

本节余下部分将利用这个命题陈述的形式性质证明埃利亚什伯格--格罗莫夫刚性定理。我们紧随埃克兰与霍费尔在 (EkH90) 中给出的证明方案。

\subsection{保容量线性映射的性质}
\label{subsec:capacity-preserving-linear-maps}

事实表明，线性映射下有界椭球与余迷向柱体的容量研究，在 (EkH90) 的刚性证明中起重要作用。

一个（开）有界椭球 $E=E_q$ 是由
\[
 E_q=\{z\in\mathbb C^n\mid q(z)<1\}
\]
定义的子集，其中 $q$ 是 $\mathbb C^n$ 上的正定二次型。作为归一化公理的一个简单推论，有如下结论。

\begin{proposition}\label{prop:capacities-linear-ellipsoid}
设 $c$ 是任意容量。则对任意有界椭球 $E$ 与任意线性映射 $L$，都有
\[
 \underline c(L\cdot E)=c(L\cdot E)=\overline c(L\cdot E).
\]
\end{proposition}

\begin{exercise}\label{ex:capacities-linear-ellipsoid}
证明这个命题。
\end{exercise}

\begin{definition}\label{def:coisotropic-cylinder}
\stfhIndex{I077}{yu^mi^xiang^zhu^ti}{余迷向柱体（coisotropic cylinder）}
设 $W\subset\mathbb C^n$ 是余迷向子空间，并设 $\Omega$ 是有界开子集。称这个和 $\Omega+W$ 为 $\Omega$ 上的\emph{余迷向柱体}。
\end{definition}

\stfhIndex{I111}{ai^ke^lan^huo^fei^er}{埃克兰--霍费尔（Ekeland–Hofer）}
注意，因为该柱体包含一个小球，所以 $c(\Omega+W)>0$。埃克兰与霍费尔 (EkH90) 证明了下面这个有趣的引理。（埃克兰与霍费尔把这个假设表述为 $W^\omega$ 是\emph{零}空间。这里改用更标准的术语\emph{“余迷向”}。）

\begin{lemma}[引理 3 (EkH90)]\label{lem:capacity-coisotropic-cylinder}
设 $W$ 是 $\mathbb C^n$ 中余维数为 $2$ 的子空间。对有界开子集
$\Omega$，考虑柱体 $\Omega+W$。则
$c(\Omega+W)=\infty$ 当且仅当 $W$ 是余迷向的。
\end{lemma}

\begin{proof}
先从下面这个简单练习开始。

\begin{exercise}\label{ex:coisotropic-codim-two-normal-form}
设 $W$ 是余维数为 $2$ 的余迷向子空间。则存在线性辛映射
$A\in Sp(2n,\mathbb R)$，使得
\[
 A\cdot W=\{z\in\mathbb C^n\mid
   \operatorname{Re}z_1=\operatorname{Re}z_2=0\}.
\]
（任意余维数的余迷向子空间集合之完整描述见 (OhP05)。）
\end{exercise}

记 $z_j=q_j+\sqrt{-1}p_j$。作一个线性辛坐标变换后，可以把
$\epsilon>0$ 取得如此之小，使得
\[
 z\in\Omega+W
 \qquad\text{只要}\qquad q_1^2+q_2^2\leq\epsilon^2.
\]
由上述练习给出的标准形，这样的 $\epsilon>0$ 存在。根据
$\Omega+W$ 的定义，对所有 $N>0$ 的椭球 $E_{q_{\epsilon,N}}$，都有
\[
 E_{q_{\epsilon,N}}\subset\Omega+W,
\]
其中二次型 $q_{\epsilon,N}$ 由
\[
 q_{\epsilon,N}
 =\left(\frac1{\epsilon^2}q_1^2+\frac1{N^2}p_1^2\right)
  +\left(\frac1{\epsilon^2}q_2^2+\frac1{N^2}p_2^2\right)
  +\frac1{N^2}\sum_{k=3}^n|z_k|^2
\]
给出。但是，在一个辛坐标变换下，这个二次型共轭于
\[
 \overline q_{\epsilon,N}
 =\frac1{\epsilon N}(|z_1|^2+|z_2|^2)
  +\frac1{N^2}\sum_{k=3}^n|z_k|^2.
\]
所以
\[
 c(E_{q_{\epsilon,N}})=c(E_{\overline q_{\epsilon,N}}).
\]
另一方面，如果 $N>\epsilon$，从而 $\epsilon N<N^2$，则
$c(E_{\overline q_{\epsilon,N}})=\pi\epsilon N$。（证明这一点！）利用
$c$ 的单调性，得到 $c(\Omega+W)\geq\pi\epsilon N$。由于这对所有
$N>0$ 都成立，已经证明 $c(\Omega+W)=\infty$。

接着证明逆命题。考虑其逆否命题，需要证明当 $W$ 不是余迷向的时，
$c(\Omega+W)<\infty$。仍从下面的练习开始。

\begin{exercise}\label{ex:noncoisotropic-codim-two-normal-form}
假设 $W$ 是一个余维数为 $2$ 且非余迷向的子空间。证明存在线性辛映射 $L$，使得
\[
 L\cdot W=\{z\in\mathbb C^n\mid z_1=0\}.
\]
\end{exercise}

由这个练习，不失一般性可假设
$W=\{z\in\mathbb C^n\mid z_1=0\}$。由于 $\Omega$ 有界，存在充分大的 $N>0$，使得
\[
 \Omega+W\subset
 \{z\in\mathbb C^n\mid |z_1|^2<N^2\}=Z^{2n}(N).
\]
于是单调性蕴含
$c(\Omega+W)\leq c(Z^{2n}(N))=\pi N^2<\infty$，证明完毕。
\end{proof}

接着证明下面的引理。

\begin{lemma}\label{lem:capacity-preserving-linear-invertible}
设 $\Psi$ 是线性映射。如果它保持任意有界椭球的一种容量 $c$，则
$\Psi$ 可逆。
\end{lemma}

\begin{proof}
假设对任意有界椭球 $S$，都有 $c(\Psi(S))=c(S)$。只需证明
$\Psi$ 满射。反设 $\Psi(\mathbb C^n)\ne\mathbb C^n$。则
$\Psi(\mathbb C^n)$ 包含在一个过原点的超曲面 $V$ 中。必要时作线性辛坐标变换，可以假设
\[
 V=\{z\in\mathbb C^n\mid\operatorname{Re}z_1=0\}.
\]
考虑任意有界椭球 $S$，例如 $S=B^{2n}(r)$。则 $\Psi(S)$ 是包含于
$V$ 的有界子集。因此存在充分大的 $N>0$，使得对所有
$\epsilon>0$，都有
\[
 \Psi(S)\subset
 \left\{z\mathrel{\bigg|}
   \frac1{\epsilon^2}q_1^2+\frac1{N^2}p_1^2
   +\frac1{N^2}\sum_{k=2}^n|z_k|^2\right\}.
\]
于是单调性蕴含对所有 $\epsilon>0$，
\[
 c(\Psi(S))\leq\pi\epsilon N,
\]
因而 $c(\Psi(S))=0$。另一方面，$0<c(S)=c(\Psi(S))$；其中不等式来自归一化公理，等式来自辛不变性。于是得到矛盾，所以 $\Psi$ 必须可逆。
\end{proof}

接着回忆线性映射的\emph{辛对偶} $\Psi^*$ 之定义。$\Psi^*$ 是满足
\[
 \omega_0(\Psi(v),w)=\omega_0(v,\Psi^*(w))
\]
的线性映射。

上述命题与引理立即给出如下推论。

\begin{corollary}\label{cor:capacity-map-coisotropic-planes}
设 $\Psi$ 是保持容量 $c$ 的线性映射。则 $\Psi$ 把任意余维数为
$2$ 的余迷向子空间映到一个余维数为 $2$ 的余迷向子空间。等价地，
$(\Psi^*)^{-1}$ 把迷向 $2$-平面映到迷向 $2$-平面。
\end{corollary}

\begin{proof}
假设 $W$ 是余维数为 $2$ 的余迷向子空间。设 $\Omega$ 是非空有界开子集，并考虑柱体 $\Omega+W$。命题 6.1.9 蕴含
$c(\Omega+W)=\infty$。由于 $\Psi$ 保持 $c$，得到
\[
 \infty=c(\Psi(\Omega+W))=c(\Psi(\Omega)+\Psi(W)).
\]
$\Psi$ 的可逆性蕴含 $\Psi(\Omega)$ 仍是非空有界开子集，而
$\Psi(W)$ 是余维数为 $2$ 的子空间。于是命题 6.1.9 蕴含
$\Psi(W)$ 是余迷向的。这证明了第一个陈述。

对于第二个陈述，设 $U$ 是维数为 $2$ 的迷向子空间，并记
$W=U^{\omega_0}$；这是余维数为 $2$ 的余迷向子空间。回忆，对任意子空间 $U\subset\mathbb C^n$ 与线性映射 $\Psi$，都有对偶公式
\[
 \Psi(U^{\omega_0})
 =\bigl((\Psi^*)^{-1}U\bigr)^{\omega_0}.
\]
如果 $U$ 是维数为 $2$ 的迷向子空间，则 $U^{\omega_0}$ 是余维数为
$2$ 的余迷向子空间，所以由第一个陈述，$\Psi(U^{\omega_0})$ 是余迷向的。于是由上述对偶公式，
$\bigl((\Psi^*)^{-1}U\bigr)^{\omega_0}$ 是余迷向的；等价地，
$(\Psi^*)^{-1}(U)$ 是迷向的。证明完毕。
\end{proof}

最后证明下面的一般引理。

\begin{lemma}\label{lem:isotropic-two-plane-conformal}
设 $(V,\omega)$ 是任意辛向量空间，而 $\Phi:V\to V$ 是把迷向
$2$-平面映到迷向 $2$-平面的可逆线性映射。则对某个 $d\ne0$，
\[
 \Phi^*\omega=d\cdot\omega.
\]
\end{lemma}

\begin{proof}
选取达布基
\[
 \{e_1,\ldots,e_n,f_1,\ldots,f_n\}
\]
使得
\[
 \omega(e_i,e_j)=0=\omega(f_i,f_j),
 \qquad \omega(e_i,f_j)=\delta_{ij}.
\]
于是，由于 $\Phi$ 保持迷向 $2$-平面，得到
\begin{equation}
 \omega(\Phi(e_i),\Phi(e_j))=0
 =\omega(\Phi(f_i),\Phi(f_j))
 \label{eq:isotropic-images-basis-pairs}
\end{equation}
以及
\begin{equation}
 \omega(\Phi(e_i),\Phi(f_j))=0
 \qquad\text{当 }i\ne j.
 \label{eq:isotropic-images-mixed-pairs}
\end{equation}
由 $\Phi$ 的可逆性和 $\omega$ 的非退化性，得到
\[
 \omega(\Phi(e_j),\Phi(f_j))=d_j
\]
对某个 $d_j\ne0$ 成立。另一方面，简单计算给出
\[
 \omega(e_i-e_j,f_i+f_j)=0;
\]
也就是说，$\{e_i-e_j,f_i+f_j\}$ 的张成空间是迷向的，因而
$\{\Phi(e_i)-\Phi(e_j),\Phi(f_i)+\Phi(f_j)\}$ 的张成空间也是迷向的。所以
\[
\begin{aligned}
 0&=\omega(\Phi(e_i)-\Phi(e_j),\Phi(f_i)+\Phi(f_j))\\
  &=\omega(\Phi(e_i),\Phi(f_i))
   -\omega(\Phi(e_j),\Phi(f_j))=d_i-d_j
\end{aligned}
\]
对所有 $i,j$ 成立。以 $d$ 表示这些 $d_j$ 的公共值。于是已经证明，对所有 $j$，
\[
 \omega(\Phi(e_j),\Phi(f_j))=d\ne0.
\]
将其与\eqref{eq:isotropic-images-basis-pairs}和
\eqref{eq:isotropic-images-mixed-pairs}结合，便证明对所有
$v,w\in\mathbb C^n$，都有
$\omega(\Phi(v),\Phi(w))=d\cdot\omega(v,w)$；即如所述，
$\Phi^*\omega=d\cdot\omega$。
\end{proof}

现在可以完成下面定理的证明。

\begin{theorem}\label{thm:linear-capacity-rigidity}
假设 $\Psi:\mathbb C^n\to\mathbb C^n$ 是保容量线性映射。则
$\Psi$ 或者是辛的，或者是反辛的。
\end{theorem}

\begin{proof}
记 $\Phi=(\Psi^*)^{-1}$。由推论 6.1.13 与引理 6.1.14，
\[
 \Phi^*\omega_0=d\cdot\omega_0.
\]
根据 $d$ 的符号，随即有
$\frac1{\sqrt{|d|}}\Phi$ 或
$\frac1{\sqrt{|d|}}\Gamma\circ\Phi$ 是辛的，其中 $\Gamma$ 是
$\mathbb C^n$ 的任意反辛同构。注意，$\Gamma$ 也保持容量。于是
\[
 \sqrt{|d|}\Psi
 =\left(\frac1{\sqrt{|d|}}\Phi^*\right)^{-1}
 \quad\text{或}\quad
 \sqrt{|d|}(\Gamma^*)^{-1}\circ\Psi.
\]
注意，$(\Gamma^*)^{-1}$ 也保持容量。所以在第一种情形有
\[
 c(\sqrt{|d|}\Psi(S))=c(S).
\]
另一方面，共形性公理蕴含
$c(\sqrt{|d|}\Psi(S))=|d|c(\Psi(S))$，故 $|d|=1$。同样的论证适用于第二种情形，因而证明了
$\Psi^*\omega_0=\pm\omega_0$。
\end{proof}

\subsection{刚性定理的证明}
\label{subsec:proof-c0-rigidity}

本小节处理一般的非线性保容量映射及其一致极限。先证明下面的定理。

\begin{theorem}\label{thm:uniform-limit-capacity-maps-derivative}
设 $\Phi_k$ 是一列从单位球 $B^{2n}(1)$ 到 $\mathbb C^n$ 的连续映射，一致收敛于 $\Phi$。假设所有 $\Phi_k$ 都保持所有有界椭球
$S\subset B^{2n}(1)$ 的容量：
\[
 c(\Phi_k(S))=c(S)\qquad\text{对所有 }k.
\]
如果 $\Phi$ 在 $0\in B^{2n}(1)$ 处可微，则导数 $d\Phi(0)$ 或者是辛的，或者是反辛的。
\end{theorem}

在给出这个定理的证明之前，先陈述它的两个直接推论。

\begin{corollary}\label{cor:uniform-limit-symplectic-embeddings-derivative}
假设 $\Phi_k:B^{2n}(1)\to\mathbb C^n$ 是一列一致收敛于映射
$\Phi:B^{2n}(1)\to\mathbb C^n$ 的辛嵌入，并且 $\Phi$ 在 $0$ 处可微。则
\[
 d\Phi(0)\in Sp(2n\mathbb R).
\]
\end{corollary}

\begin{proof}
由于辛嵌入保持任意容量，定理 6.1.16 蕴含 $d\Psi(0)$ 或者是辛的，或者是反辛的。还需证明 $d\Psi(0)$ 确实是辛的。

注意，所有 $\Phi_k$ 都保定向。令 $y_k=\Phi_k(0)$，并考虑映射
\[
 x\longmapsto\widetilde\Phi_k(x):=\Phi_k(x)-y_k.
\]
由于 $\Phi_k$ 一致收敛，当 $k\to\infty$ 时，
$y_k\to y_\infty$。所以映射 $\widetilde\Phi_k$ 一致收敛于
\[
 \widetilde\Phi(x):=\Phi(x)-y_\infty.
\]
显然 $d\Phi(0)=d\widetilde\Phi(0)$，因而可以改为考虑
$\widetilde\Phi_k$。这使我们不失一般性地假设 $\Phi_k(0)=0$；下文都作此假设。

考虑一个小球 $B^{2n}(\delta)$。由于 $\Phi_k(0)=0$ 且
$\Phi_k$ 保定向，环绕数为
\[
 \deg(\Phi_k|_{\partial B^{2n}(\delta)};0)=1.
\]
因为环绕数在一致极限下不变，还有
\[
 \deg(\Phi|_{\partial B^{2n}(\delta)};0)=1.
\]
由于假设 $\Phi$ 在 $0$ 处可微，
\[
 \deg(\Phi|_{\partial B^{2n}(\delta)};0)
 =\operatorname{sign}\det d\Phi(0).
\]
所以已经证明 $\operatorname{sign}\det d\Phi(0)=1$，从而
$d\Phi(0)$ 保定向。

此时，如果 $n$ 是奇数，则 $d\Phi(0)$ 必为辛的，因为反辛映射在这些维数中反定向。当 $n$ 是偶数时，考虑定义域为
$\mathbb C\times B^{2n}(1)$ 的乘积嵌入序列
$\operatorname{Id}_{\mathbb C}\times\Phi_k$。它收敛于
$\operatorname{Id}_{\mathbb C}\times\Phi$。显然，每个
$\operatorname{Id}_{\mathbb C}\times\Phi_k$ 都定义了
$B^{2(n+1)}(1)$ 的一个辛嵌入，而 $n+1$ 为奇数。此外，极限
$\operatorname{Id}_{\mathbb C}\times\Phi$ 在 $0$ 处的导数为
$\operatorname{Id}_{\mathbb C}\times d\Phi(0)$。所以
$\operatorname{Id}_{\mathbb C}\times d\Phi(0)$ 是辛的，进而
$d\Phi(0)$ 也是辛的。
\end{proof}

这个推论立即给出我们想证明的刚性结果。

\begin{theorem}\label{thm:symp-closed-compact-open}
设 $(M,\omega)$ 是任意辛流形。则 $\Symp(M,\omega)$ 在
$\Diff(M)$ 的紧开 $C^0$ 拓扑中是闭的。
\end{theorem}

\begin{proof}
利用达布坐标卡作局部化，再应用推论 6.1.17。
\end{proof}

现在给出定理 6.1.16 的证明。

\begin{proof}[定理 6.1.16 的证明]
我们将证明：定理中关于任意容量 $c$ 的假设，对任意有界椭球 $S$ 都给出恒等式
\begin{equation}
 \overline c(d\Phi(0)(S))=\overline c(S).
 \label{eq:derivative-preserves-upper-capacity}
\end{equation}
然后对 $d\Phi(0)$ 应用线性刚性结果，即可完成证明。所以只需证明
\eqref{eq:derivative-preserves-upper-capacity}。

设 $\overline S$ 是满足 $\overline S\subset B^{2n}(1)$ 的有界椭球。由
$\overline c$ 的定义，存在辛微分同胚 $\Psi$，使得
\[
 \Psi^{-1}(\Phi(S))\longrightarrow Z^{2n}(r+\epsilon),
\]
其中 $\pi r^2=\overline c(\Phi(S))$。于是可以找到包含
$\Psi^{-1}(\Phi(S))$ 的有界椭球
$\widetilde S\subset Z^{2n}(r+\epsilon)$。所以
$\Psi(\widetilde S)\supset\Phi(S)$，并且由 $\overline c$ 的单调性以及归一化
$\overline c(Z^{2n}(r+\epsilon))=\pi(r+\epsilon)^2$，有
\[
 \overline c(\Psi(\widetilde S))\leq\pi(r+\epsilon)^2.
\]
由于 $\Phi_k\to\Phi$ 在 $B^{2n}(1)\supset S$ 上一致收敛，所以对所有充分大的 $k$，
\[
 \Phi_k(S)\subset\Psi(\widetilde S).
\]
$\overline c$ 的单调性与\eqref{eq:capacities-agree-on-balls}给出
\[
\begin{aligned}
 \overline c(S)=c(S)=c(\Phi_k(S))
 &\leq c(\Phi_k(S))
 \leq c(\Psi(\widetilde S))\\
 &=\overline c(\widetilde S)
 \leq\pi(r+\epsilon)^2.
\end{aligned}
\]
因为这对所有 $\epsilon>0$ 都成立，已经证明对所有有界椭球都有
\[
 \overline c(S)\leq\pi r^2=\overline c(\Phi(S)).
\]

现在考虑 $t\in(0,1)$ 时的亚历山大同痕
$\Phi_t:B^{2n}(1)\to\mathbb C^n$，定义为
\[
 \Phi_t(x)=\frac1t\Phi(tx).
\]
注意，$t\cdot S$ 也是包含于 $B^{2n}(1)$ 的有界椭球，因而
\[
 \overline c(tS)\leq\overline c(\Phi(tS)).
\]
除以 $t$，得到
\begin{equation}
 \overline c(S)
 \leq\frac1t\overline c(\Phi(tS))
 =\overline c\left(\frac1t\Phi(tS)\right)
 =\overline c(\Phi_t(S))
 \label{eq:alexander-isotopy-capacity-bound}
\end{equation}
对所有 $t\in(0,1)$ 成立。

\begin{lemma}\label{lem:capacity-alexander-limit}
有
\[
 \lim_{t\to0}\overline c(\Phi_t(S))
 \leq\overline c(d\Phi(0)(S)).
\]
\end{lemma}

\begin{proof}
注意，对每个固定的 $t\in(0,1)$，当 $k\to\infty$ 时，映射
$\Phi_{k,t}\to\Phi_t$ 在 $B^{2n}(1)$ 上一致收敛。选取序列
$t_k\to0$，并考虑对角子序列 $\Phi_{k,t_k}$。这个序列一致收敛于
$d\Phi(0)$，并且保持有界椭球的容量 $c$。所以可以应用与
\eqref{eq:alexander-isotopy-capacity-bound}证明中相同的论证，得到
\[
 \overline c(\Phi_{k,t_k}(S))
 \leq\overline c(d\Phi(0)(S)).
\]
另一方面，由于 $\Phi$ 在 $0$ 处可微，$\Phi_t$ 在 $B^{2n}(1)$ 上一致收敛于 $d\Phi(0)$。由于 $\Phi_{t_k}$ 与 $\Phi_{k,t_k}$ 都一致收敛于
$d\Phi(0)$，随即有
\[
 \lim_{k\to\infty}\overline c(\Phi_{t_k})
 =\lim_{k\to\infty}\overline c(\Phi_{k,t_k}).
\]
这个收敛性与上述不等式完成证明。
\end{proof}

把这个引理与\eqref{eq:alexander-isotopy-capacity-bound}结合，已经证明
\begin{equation}
 \overline c(S)\leq\overline c(d\Phi(0)(S)).
 \label{eq:derivative-capacity-lower-bound}
\end{equation}
由与引理 6.1.12 证明中相同的论证，从这里推出 $d\Phi(0)$ 可逆。

现在只需证明相反的不等式。给定 $\gamma>0$，并考虑满足
\[
 (1+\gamma)\tau<1
\]
的 $\tau>0$。

\begin{lemma}\label{lem:uniform-map-covers-derivative-image}
对充分小的 $\tau>0$，所有充分大的 $k$ 都满足
\[
 \Phi_k((1+\gamma)\tau S)\supset d\Phi(0)(\tau S).
\]
\end{lemma}

\begin{proof}
先回忆 $\Phi_k(0)=0=\Phi(0)$。由于 $\Phi$ 在 $0$ 处可微，存在连续函数
$\epsilon:(0,1)\to(0,\infty)$，使得当 $s\to0$ 时
$\epsilon(s)\to0$，并且
\[
 |\Phi(x)-d\Phi(0)x|\leq\epsilon(|x|)|x|.
\]
此外，由于 $\Phi_k\to\Phi$ 一致收敛，对任意给定的
$\delta\in(0,1)$，存在 $k(\delta)$，使得对所有 $k\geq k(\delta)$，
\begin{equation}
 |\Phi_k(x)-d\Phi(0)x|
 \leq\epsilon(|x|)|x|+\delta.
 \label{eq:uniform-derivative-error}
\end{equation}
令
\[
 \delta=\epsilon((1+\gamma)\tau)(1+\gamma)\tau.
\]
考虑 $d\Phi(0)$ 与 $\Phi_k$ 之间的同伦
\[
 t\in[0,1]\longmapsto
 h_t:=d\Phi(0)+t(\Phi_k-d\Phi(0)).
\]
因为 $d\Phi(0)$ 可逆，存在 $d>0$，使得对所有 $x$，
\[
 |d\Phi(0)x|\geq d|x|.
\]
另一方面，只要 $\tau>0$ 充分小，由
\eqref{eq:uniform-derivative-error}可得，对所有 $x\in S$，
\[
\begin{aligned}
 &|\Phi_k((1+\gamma)\tau x)
   -d\Phi(0)(1+\gamma)\tau x|\\
 &\qquad\leq
 \epsilon((1+\gamma)\tau|x|)(1+\gamma)\tau|x|
 +\epsilon((1+\gamma)\tau)(1+\gamma)\tau.
\end{aligned}
\]
所以，如果 $x\in\partial S\subset B^{2n}(1)$ 且
$y\in d\Phi(0)(\tau S)$，并且对某个 $z\in S$ 有 $y=\tau z$，则
\[
\begin{aligned}
 &|h_t((1+\gamma)\tau x)-d\Phi(0)(\tau z)|\\
 &\quad={}
 |d\Phi(0)((1+\gamma)\tau x-\tau z)
 +t(\Phi_k((1+\gamma)\tau x)
       -d\Phi(0)((1+\gamma)\tau x))|\\
 &\quad\geq d\tau|(1+\gamma)x-z|
       -2\epsilon((1+\gamma)\tau)(1+\gamma)\tau\\
 &\quad\geq\tau\bigl(
   dd_1(S,\gamma)-2\epsilon((1+\gamma)\tau)(1+\gamma)\bigr),
\end{aligned}
\]
其中 $d_1(S,\gamma)$ 是一个只依赖于 $S$ 的形状与 $\gamma$ 的常数。只要 $\tau$ 小到满足
\[
 dd_1(S,\gamma)>2\epsilon((1+\gamma)\tau)(1+\gamma),
\]
最后一项就大于零。因此
\begin{equation}
 h_t(\partial(1+\gamma)\tau S)
 \cap d\Phi(0)\tau S=\varnothing.
 \label{eq:degree-homotopy-avoids-image}
\end{equation}

我们想证明 $y\in\Phi_k((1+\gamma)\tau S)$。为此，当 $t$ 与
$z\in S$ 变化时，考虑布劳威尔度数
\[
 \deg(h_t|_{\partial(1+\gamma)\tau S};y).
\]
由\eqref{eq:degree-homotopy-avoids-image}，这个度数有定义，并且在
$(t,z)\in[0,1]\times S$ 上连续。但是，当 $t=0$ 时，
$h_0=d\Phi(0)$ 是线性映射，因而
$\deg(h_0|_{\partial(1+\gamma)\tau S};0)=1$。由于
$h_1=\Phi_k$，由度数的同伦不变性，
\[
 \deg(\Phi_k|_{\partial(1+\gamma)\tau S};y)=1
\]
对所有 $y\in S$ 成立。这就完成证明。
\end{proof}

于是，容量的单调性与辛不变性、$\overline c\geq c$，以及
$d\Phi(0)$ 是线性映射这一事实蕴含
\[
\begin{aligned}
 \overline c((1+\gamma)\tau S)
 &=\overline c(\Phi_k((1+\gamma)\tau S))
 \geq c(\Phi_k((1+\gamma)\tau S))\\
 &\geq c(d\Phi(0)(\tau S))
 =\overline c(d\Phi(0)(\tau S)).
\end{aligned}
\]
这蕴含
\[
 \overline c(d\Phi(0)(S))
 \leq(1+\gamma)^2\overline c(S)
\]
对所有 $\gamma>0$ 成立。因此得到
\[
 \overline c(d\Phi(0)(S))\leq\overline c(S).
\]
将其与\eqref{eq:derivative-capacity-lower-bound}结合，便证明
\[
 \overline c(d\Phi(0)(S))=\overline c(S)
\]
对所有有界椭球 $S$ 成立。于是定理 6.1.15 蕴含 $d\Phi(0)$ 或者是辛的，或者是反辛的。这就完成了定理 6.1.16 的证明。
\end{proof}
